\section{Conclusion}

This work demonstrates the feasibility of adapting the MeanFlow one-step generative modeling for image denoising. By replacing the prior distribution with noisy images, the model effectively transforms corrupted inputs into clean outputs, achieving substantial PSNR gains with minimal function evaluations. Experiments on MNIST show that a single evaluation provides significant noise reduction, while additional steps over 2-NFE offer few returns. This success relies on the smoothness of Gaussian noise and the capacity of the DiT to capture global dependencies and preserve image structure. Limitations include restriction to MNIST, focuses beyond Gaussian noise, computational overhead from JVP and fixed CFG settings. Future work should explore more complex datasets, non-Gaussian noise, efficient training targets, and adaptive guidance strategies. Overall, this study confirms that MeanFlow can serve as an efficient, one-step denoising approach and lays groundwork for broader image restoration applications
